\section{Module \texttt{common.c}}

Ce module constitue l'épine dorsale visuelle du projet. Au lieu de laisser le développeur réécrire les mêmes blocs de code pour chaque widget (marges, expansion, alignement), \texttt{common.c} factorise toutes les propriétés héritées de la classe \texttt{GtkWidget} de GTK. Cela permet de garantir que l'application d'un style est unifiée et prédictible, quel que soit le composant instancié.

\subsection{Fonction \texttt{apply\_widget\_style}}

Cette fonction agit comme un intercepteur universel. Au lieu d'appeler manuellement \texttt{gtk\_widget\_set\_margin\_*} pour chaque bouton ou label, le développeur passe une sous-structure \texttt{style}. L'implémentation est conçue de manière défensive : elle vérifie méthodiquement la présence de pointeurs non nuls et évalue des drapeaux booléens (comme \texttt{set\_halign}) pour garantir qu'aucune propriété GTK n'est écrasée involontairement si elle n'a pas été explicitement demandée par l'utilisateur.

\begin{lstlisting}[language=C]
void apply_widget_style(GtkWidget *widget, const widget_style_config *style) {
    if (widget == NULL || style == NULL) {
        return;
    }

    if (style->css_class && style->css_class[0] != '\0') {
        gtk_widget_add_css_class(widget, style->css_class);
    }
    if (style->css_classes != NULL) {
        for (int i = 0; style->css_classes[i] != NULL; i++) {
            gtk_widget_add_css_class(widget, style->css_classes[i]);
        }
    }

    if (style->width_request != 0 || style->height_request != 0) {
        gtk_widget_set_size_request(widget, style->width_request, style->height_request);
    }

    gtk_widget_set_margin_top(widget, style->margin_top);
    gtk_widget_set_margin_bottom(widget, style->margin_bottom);
    gtk_widget_set_margin_start(widget, style->margin_start);
    gtk_widget_set_margin_end(widget, style->margin_end);

    if (style->set_halign) {
        gtk_widget_set_halign(widget, style->halign);
    }
    if (style->set_valign) {
        gtk_widget_set_valign(widget, style->valign);
    }
    if (style->set_hexpand) {
        gtk_widget_set_hexpand(widget, style->hexpand);
    }
    if (style->set_vexpand) {
        gtk_widget_set_vexpand(widget, style->vexpand);
    }
    if (style->set_sensitive) {
        gtk_widget_set_sensitive(widget, style->sensitive);
    }
    if (style->set_visible) {
        gtk_widget_set_visible(widget, style->visible);
    }

    if (style->tooltip && style->tooltip[0] != '\0') {
        gtk_widget_set_tooltip_text(widget, style->tooltip);
    }
}
\end{lstlisting}

\paragraph{Macros associées :}
L'implémentation de la configuration visuelle est complétée par des macros de substitution. La première série mappe nos propres constantes (ex: \texttt{ALIGN\_CENTER}) directement sur les types énumérés internes de GTK. Cette approche encapsule le code utilisateur en le découplant légèrement de l'API sous-jacente tout en conservant une résolution à la compilation, assurant une performance maximale sans coût d'exécution. La seconde série (\texttt{MARGIN\_*} ) instaure un système de design (Design System) par paliers fixes. En forçant l'utilisation de ces constantes plutôt que des valeurs magiques éparpillées (magic numbers), le wrapper garantit mécaniquement une harmonie spatiale et une cohérence visuelle dans l'ensemble de l'application.

\begin{lstlisting}[language=C]
#define ALIGN_START  GTK_ALIGN_START
#define ALIGN_CENTER GTK_ALIGN_CENTER
#define ALIGN_END    GTK_ALIGN_END
#define ALIGN_FILL   GTK_ALIGN_FILL

#define MARGIN_TINY   4
#define MARGIN_SMALL  8
#define MARGIN_MEDIUM 16
#define MARGIN_LARGE  24
\end{lstlisting}

\subsection{APIs GTK utilisées}

\begin{itemize}
    \item \textbf{\texttt{gtk\_widget\_add\_css\_class()}} : Ajoute une classe de style CSS spécifiée au widget cible. Dans \texttt{apply\_widget\_style}, elle permet d'appliquer de manière dynamique une ou plusieurs classes CSS définies dans la configuration pour modifier l'apparence visuelle du widget.
    \item \textbf{\texttt{gtk\_widget\_set\_size\_request()}} : Fixe les dimensions minimales (largeur et hauteur) que le widget doit tenter d'occuper. Le wrapper utilise cette fonction pour forcer une taille spécifique demandée par le développeur dans la structure \texttt{widget\_style\_config}.
    \item \textbf{\texttt{gtk\_widget\_set\_margin\_top()}} (et variantes) : Définit les marges extérieures du widget pour chaque direction. Le module les exploite pour assurer un espacement cohérent entre les widgets sans avoir à manipuler directement les conteneurs parents.
    \item \textbf{\texttt{gtk\_widget\_set\_halign()}} et \textbf{\texttt{gtk\_widget\_set\_valign()}} : Contrôlent l'alignement du widget respectivement sur les axes horizontal et vertical. Elles permettent au wrapper de gérer la position du widget (centré, aligné à gauche/droite, ou étendu) dans sa zone d'affichage.
    \item \textbf{\texttt{gtk\_widget\_set\_hexpand()}} et \textbf{\texttt{gtk\_widget\_set\_vexpand()}} : Indiquent au gestionnaire de mise en page si le widget doit absorber l'espace supplémentaire disponible. Ces appels sont cruciaux pour définir le comportement de redimensionnement réactif des interfaces construites avec le wrapper.
    \item \textbf{\texttt{gtk\_widget\_set\_tooltip\_text()}} : Associe un texte d'aide contextuelle qui apparaît lorsque l'utilisateur survole le widget avec la souris. Elle est systématiquement appelée par le module si une chaîne de caractères \texttt{tooltip} est fournie.
\end{itemize}

\newpage
